\exercisesection{9}

\begin{enumerate}
\def\labelenumi{\arabic{enumi}.}
\tightlist
\item
  \begin{enumerate}
  \def\labelenumii{(\alph{enumii})}
  \tightlist
  \item
    Let K be a field (not necessarily commutative) and T a locally compact topology on K compatible with the ring structure on K; show that T is compatible with the field structure on K. (Use the Theorem of R. Ellis (General Topology, Chapter X, § 3, Exercise 25) or argue directly reproducing the proofs of Integration, Chapter VII, § 1, no. 10 and those of Proposition 1 of this paragraph.)
  \end{enumerate}
\end{enumerate}

\begin{enumerate}
\def\labelenumi{(\alph{enumi})}
\setcounter{enumi}{1}
\tightlist
\item
  Give an example of a locally compact topology on a commutative field K, which is compatible with the additive group structure on K, but not with its ring structure. (Take K to be the field of fractions of a compact integral domain A (A being for example a ring of formal power series \(k[[X, Y]]\) , where k is finite) and take as fundamental system of neighbourhoods of 0 in K the neighbourhoods of 0 in A.)
\end{enumerate}

¶ 2. (a) In a space \(\mathbf{R}^n\) ( \(n \geqslant 2\) ), let U be a non-empty open set with non-empty complement; show that the frontier of U contains a non-empty perfect set (cf.~General Topology, Chapter I, § 9, Exercise 17).

\begin{enumerate}
\def\labelenumi{(\alph{enumi})}
\setcounter{enumi}{1}
\item
  Deduce from (a) that, if \(\mathbf{A}\) is an everywhere dense subset of \(\mathbf{R}^n\) which meets every perfect set in \(\mathbf{R}^n\) , \(\mathbf{A}\) is convex.
\item
  Show that there exists in \(\mathbf{C}\) an everywhere dense subfield \(\mathbf{K}\) which is connected and locally connected and is a pure transcendental extension of \(\mathbf{Q}\) (apply General Topology, Chapter IX, § 5, Exercise 18 (b) and (c), constructing \(\mathbf{K}\) by transfinite induction, using (b) and the method described in Set Theory, Chapter III, § 6, Exercise 24). Deduce that there exists a subfield \(\mathbf{K}' \supset \mathbf{K}\) of \(\mathbf{C}\) which is (algebraically) isomorphic to \(\mathbf{R}\) , connected and locally connected.
\item
  Show, using (c), that there exists on C the topology of a connected and locally connected space, which is compatible with the field structure and for which the completion of C is an algebra over C which is the direct composition of two fields isomorphic to C.
\end{enumerate}

\begin{enumerate}
\def\labelenumi{\arabic{enumi}.}
\setcounter{enumi}{2}
\tightlist
\item
  Let K be a totally disconnected non-discrete locally compact commutative field; let A be the ring of the absolute value mod \(_{K}\) on K and let U be the group of units of A. For every integer n \textgreater{} 0, let \(_{n}U\) denote the group of n-th roots of unity in K and \(U^{n}\) the subgroup of U consisting of the n-th powers of the elements of U. Show that, if n is prime to the characteristic of K, \(U^{n}\) is an open subgroup of U and that
\end{enumerate}

\[
\operatorname{Card} (\mathrm{U} / \mathrm{U} ^ {n}) = \operatorname{mod} _ {\mathrm{K}} (n). \operatorname{Card} (_ {n} \mathrm{U})
\]

(use Exercise 14 of Integration, Chapter VII, §2 and Commutative Algebra, Chapter III, §4, no. 6, Corollary 1 to Theorem 2 to show that, if m is the maximal ideal of A, the image under \(x \mapsto x^n\) of \(1 + \mathfrak{m}^k\) is \(1 + n \cdot \mathfrak{m}^k\) for \(k\) sufficiently large).

\begin{enumerate}
\def\labelenumi{\arabic{enumi}.}
\setcounter{enumi}{3}
\tightlist
\item
  \begin{enumerate}
  \def\labelenumii{(\alph{enumii})}
  \tightlist
  \item
    Let K be a commutative field and v a valuation on K such that K is Henselian with v (§ 8, Exercise 6). Suppose further that K and the residue field k of v are of characteristic 0. Show that there exists a subfield \(K_{0}\) of the ring A of v such that the canonical mapping \(A \rightarrow k\) , restricted to \(K_{0}\) , is an isomorphism of \(K_{0}\) onto k. (Let H be a subfield of A such that the image of H under the canonical mapping is an isomorphism onto a subfield E of k; show that, if \(E \neq k\) , there exists \(\alpha \notin H\) in A such that the subfield \(\mathrm{H}(\alpha)\) of K is contained in A and is canonically isomorphic to \(\mathrm{E}(\bar{\alpha})\) , where \(\bar{\alpha}\) is the class of \(\alpha\) in k; distinguish two cases according to whether \(\bar{\alpha}\) is algebraic or transcendental over E.)
  \end{enumerate}
\end{enumerate}

\begin{enumerate}
\def\labelenumi{(\alph{enumi})}
\setcounter{enumi}{1}
\tightlist
\item
  Suppose further that v is a discrete valuation and that K is complete with respect to v; deduce from (a) that K is isomorphic to the field of formal power series \(k((\mathbf{T}))\) .
\end{enumerate}

\begin{enumerate}
\def\labelenumi{\arabic{enumi}.}
\setcounter{enumi}{4}
\tightlist
\item
  \begin{enumerate}
  \def\labelenumii{(\alph{enumii})}
  \tightlist
  \item
    Suppose that K is a commutative field, v a valuation of height 1 on K such that K is complete with respect to v and A the ring of v; suppose further that the residue field k of v is perfect and of characteristic p \textgreater{} 0. For every element \(\xi \in k^{*}\) and every integer n, let \(x_{n}\) be an element of the class \(\xi^{p-n}\) in A; show that the sequence \((x_{n}^{p^{n}})\) is a Cauchy sequence in A whose limit is independent of the choice of the \(x_{n}\) in the classes \(\xi^{p-n}\) . If this limit is denoted by \(\phi(\xi)\) , show that \(\phi\) is the unique isomorphism u of the multiplicative group \(k^{*}\) to the multiplicative group \(K^{*}\) such that, for all \(\xi \in k^{*}\) , \(u(\xi)\) is an element of A belonging to the class \(\xi\) .
  \end{enumerate}
\end{enumerate}

\begin{enumerate}
\def\labelenumi{(\alph{enumi})}
\setcounter{enumi}{1}
\item
  If \(\mathbf{K}\) is also of characteristic \(p\) , show that \(\phi\) , extended to \(k\) by \(\phi(0) = 0\) , is an isomorphism of the field \(k\) onto a subfield of \(\mathbf{K}\) . Deduce a new proof of Theorem 1 (iii) of no. 3.
\item
  Suppose that \(k\) is finite. Show that, if \(r\) is prime to \(p\) , the group \((\mathbf{K}^*)^r\) of \(r\) -th powers of elements of \(\mathbf{K}^*\) is of finite index in \(\mathbf{K}^*\) (use Hensel's Lemma). Show that, if further \(v\) is discrete and \(\mathbf{K}\) of characteristic 0, the same result is valid without restriction on \(r\) (observe that every element of \(1 + p^2\mathbf{A}\) is a \(p\) -th power).
\end{enumerate}

